\begin{definition}
Let $\Omega\subset\mathbb{C}$ be open, let $n\ge1$ be an integer and let
$1<p<\infty$.  A mapping $H\colon\mathbb{C}\to\mathbb{C}$ is said to
\emph{satisfy the system (18) on $\Omega$} if there are weak Wirtinger
derivatives $a=\partial H/\partial\xi$ and
$b=\partial H/\partial\bar\xi$ of $H$ on $\Omega$, in the sense of
\noderef{IsWeakWirtingerOn}, such that
\[
  b(\xi)
  =\Bigl(\frac12-\frac1p\Bigr)
   \Bigl[\frac{\xi}{\bar\xi}\,a(\xi)
     +\frac{\bar\xi^{\,n}}{\xi^{\,n}}\,\overline{a(\xi)}\Bigr]
\]
at almost every $\xi\in\Omega$ with $\xi\neq0$, ``almost every'' referring to
two-dimensional Lebesgue measure.

This is equation (18) of the paper, the linear first order system that the
hodograph transformation of Section~3 produces for the inverse $H$ of the
quasiconformal factor $\chi$, written
\[
  H_{\bar\xi}
  =\Bigl(\frac12-\frac1p\Bigr)
   \Bigl[\frac{\xi}{\bar\xi}H_\xi+\frac{\bar\xi^{\,n}}{\xi^{\,n}}\,\overline{H_\xi}\Bigr].
  \tag{18}
\]
Its coefficients $\xi/\bar\xi$ and $\bar\xi^{\,n}/\xi^{\,n}$ are defined only
for $\xi\neq0$, which is why the origin is excluded; the single point $\xi=0$
has measure zero.

The derivatives are the weak ones, for the same reason as in
\noderef{QuasiregularOn}: the only hypothesis the paper's solutions of (18)
carry is membership in $W^{1,2}$, so $H_\xi$ and $H_{\bar\xi}$ are the
distributional derivatives, and a solution in this sense is a condition on the
class of $H$ rather than on a chosen representative.  Read through the real
Fr\'echet derivative of the given representative the system would be satisfied
by every nowhere differentiable representative of every $L^2$-class, both
sides vanishing identically, and the paper's assertion that a $W^{1,2}$
solution of (18) is $K$-quasiregular with $K=\max\{p-1,1/(p-1)\}$ would be
false.  By the next-to-last paragraph of \noderef{IsWeakWirtingerOn}, whenever
$H$ is in addition differentiable at almost every point of $\Omega$ the pair
$(a,b)$ agrees almost everywhere with the pointwise pair
$(H_\xi,H_{\bar\xi})$ read through $\mathrm{d}H$; this is how Sections~4--6,
which work with almost everywhere differentiable maps, recover the pointwise
form of (18).  By the last paragraph of \noderef{IsWeakWirtingerOn} the choice
of the witnessing pair is immaterial.
\end{definition}
